% $Id: vectors.tex 1784 2012-04-27 23:29:31Z  $
% !TEX root = main.tex

\chapter{Replication Backbone and Causal Consistency}\label{chap:backbone}

\section{Introduction}

The \textit{CRDTBackbone} module realises a generic operation based replication protocol.  
It alligns the vector clock definitions from Chapter~\ref{chap:vectors} with the payload semantics from Chapter~\ref{chap:payload}.  
By associating each update with a timestamp and enforcing causal delivery.  
This chapter explains the organisation of the backbone, illustrates its main algorithms and sketches why they are correct.
\textbf{TO DO: Add a sentence explaining how the generic operation can be used for other communication/consistency models.}

\section{Configurations and vector clocks}

The backbone maintains, for each replica, its local payload state, a vector clock summarising the set of updates it has seen, and a network buffer of pending messages.  
Concretely, a configuration comprises three components:
\begin{itemize}
  \item A finite map \(\mathit{states}:\mathsf{Ide}\to S\) mapping each replica identifier to its current payload state (with \(S\) the payload’s state domain);
  \item \textbf{TO DO:a map?}A \emph{registry} \(\mathit{clocks}:\mathsf{Ide}\to VC\) storing, for each replica\textbf{TO DO: Should S be parametrized by the map?}, its current vector timestamp \cite{lamport1978time,fidge1988timestamps,mattern1989virtual};
  \item A multiset \(\mathit{net}\) of messages, each message being a triple \((\text{destination},\text{payload},\text{timestamp})\).
  \textbf{TO DO:write in theory sets notation. net = \{(\text{destination},\text{payload},\text{timestamp}) | destination \in Ide, payload \in V, timestamp \in VC n:Ints is the multiplicity of the message
  \}
  
  add a note on why the messages need to be multisets } 

\end{itemize}

\subsection{Configuration type and basic accessors}

In Athena we model a configuration using a record with three fields: a map of actor states \textbf{TO DO:astates}, a registry of vector clocks \textbf{TO DO:clock's}, and a multiset of messages.  
The following snippet \textbf{TO DO:ref/snippets} shows the type definition and simple accessor functions:

\AthenaCode{lst:backbone-config-type}{Configuration datatype}
\begin{lstlisting}[language=Athena,style=athena]
datatype (Configuration State Value) :=
  (kfg astates:(FMap.Map Ide State)
       clocks:Registry
       net:(MSet.MSet (Message Value Clock)))
\end{lstlisting}
\textbf{TO DO: this doesn't correspond to the previous definition, adjust.}

\AthenaCode{lst:backbone-config-access}{Configuration accessors}
\begin{lstlisting}[language=Athena,style=athena]
define (get-cfg-states cfg) :=
  match cfg {
    (Backbone.kfg states _ _) => states
  | _ => empty-map 
  }

define (get-cfg-clocks cfg) :=
  match cfg {
    (Backbone.kfg _ clks _) => clks
  | _ => registry-empty 
  }

define (get-cfg-net cfg) :=
  match cfg {
    (Backbone.kfg _ _ network) => network
  | _ => empty-mset 
  }
\end{lstlisting}

The record constructor \lstinline{Backbone.kfg} encapsulates three components together.  
The helper functions \lstinline{get-cfg-states}, \lstinline{get-cfg-clocks} and \lstinline{get-cfg-net} simply select the appropriate field.  
These accessors are used throughout the backbone to inspect the configuration without showing its complete structure.

\section{The \texttt{CRDTBackbone} module}

We introduce the core functions of the CRDT Backbone module in four groups: local clock management, message delivery, broadcasting, and initialisation.  
For clarity we omit low‑level helper functions such as finite map constructors and list utilities, focusing instead on the semantics.

\subsection{Local ticks and vector clocks}

When a replica performs a local update, it must advance its vector clock.  
The function \lstinline{cfg-tick} returns a new configuration with the caller’s clock incremented accordingly, leaving the state and network unchanged.  
The helper \lstinline{cfg-actor-clock} retrieves the current clock of a given actor.  
Both functions are defined by pattern matching on the configuration record:

\AthenaCode{lst:backbone-tick}{Updating and querying vector clocks}
\begin{lstlisting}[language=Athena,style=athena]
define (cfg-tick cfg actor) :=
  (Backbone.kfg (get-cfg-states cfg)
                (registry-tick (get-cfg-clocks cfg) actor)
                (get-cfg-net cfg))

define (cfg-actor-clock cfg actor) :=
  (let ((result:Clock (registry-get (get-cfg-clocks cfg) actor)))
    result)
\end{lstlisting}

Invoking \lstinline{cfg-tick cfg actor} produces a new configuration in which only the vector clock registry is updated:  
the entry for \lstinline{actor} is incremented using \lstinline{registry-tick}.  
This corresponds to the local event of generating a new update.  
The function \lstinline{cfg-actor-clock} simply looks up the replica’s clock in the registry; it is used when constructing messages to find the correct timestamp.

\subsection{Message delivery}



The \texttt{reduce} function models the delivery of a single message.  
If the message queue is empty, the configuration is returned unchanged.  
Otherwise, one message \lstinline{(msg id payload clock)} is removed from the multiset.  
The recipient’s current state is obtained via \lstinline{cfg-state-lookup};  
the payload update is applied using \lstinline{payload-apply};  
and the recipient’s clock registry is updated by merging in the incoming clock via \lstinline{registry-merge}.  
A new configuration with the updated state map, clock registry and reduced message network is returned.  
The helper \lstinline{reduce*} simply applies \lstinline{reduce} until no messages remain.
\textbf{TO DO: reference the snippet in the text.}  

\AthenaCode{lst:backbone-reduce}{Message delivery and repetition}
\begin{lstlisting}[language=Athena,style=athena]
define (reduce ops cfg) :=
  match (get-cfg-net cfg) {
    MSet.null => cfg
  | (MSet.insert (Payload.msg id payload clock) rest) =>
      (let ((states (get-cfg-states cfg))
            (clocks (get-cfg-clocks cfg))
            (current (cfg-state-lookup states id (payload-init ops)))
            (updated-state (payload-apply ops payload current))
            (next-states (FMap.update (pair id updated-state) states))
            (next-clocks (clock-merge-entry clocks id clock)))
        (Backbone.kfg next-states next-clocks rest))
  }

define (reduce* ops cfg) :=
  match (get-cfg-net cfg) {
    MSet.null => cfg
  | (MSet.insert _ _) => (reduce* ops (reduce ops cfg))
  }
\end{lstlisting}

The check \lstinline{cfg-state-lookup states id (payload-init ops)} returns a default initial state if the recipient has not yet been initialised;  
this makes the backbone robust to late joins.  
After applying the payload, the vector clock registry is updated via \lstinline{clock-merge-entry}, which calls \lstinline{registry-merge} to compute the componentwise maximum.  
Subsequent messages carrying smaller or equal timestamps will be discarded by the causal ordering checks in \lstinline{reduce}.  
The \lstinline{reduce*} function is used to reduce a configuration to a quiescent state where no messages remain.

\subsection{Broadcasting updates}

\textbf{TO DO: the function name must be referenced properly in text.}

When a replica performs a local update, the backbone must propagate this operation to all other replicas.  
The helper \lstinline{enqueue-message} inserts a message into the network buffer, and \lstinline{enqueue-over-states} iterates over all actors, enqueuing a message for each.  
The function \lstinline{broadcast-from} combines these primitives.

\AthenaCode{lst:backbone-enqueue}{Enqueueing updates}
\begin{lstlisting}[language=Athena,style=athena]
define (enqueue-message cfg recipient payload clock:Clock) :=
  (Backbone.kfg (get-cfg-states cfg)
                (get-cfg-clocks cfg)
                (MSet.insert (Payload.msg recipient payload clock) 
                             (get-cfg-net cfg)))

define (enqueue-over-states cfg states payload clock-fn) :=
  match states {
    FMap.empty-map => cfg
  | (FMap.update (pair actor _) rest) =>
      (let ((next (enqueue-message cfg actor payload (clock-fn actor))))
        (enqueue-over-states next rest payload clock-fn))
  | _ => cfg 
  }
\end{lstlisting}

\AthenaCode{lst:backbone-broadcast-funcs}{Broadcast functions}
\begin{lstlisting}[language=Athena,style=athena]
define (broadcast-from cfg actor payload) :=
  (let ((advanced (cfg-tick cfg actor)))
    (let ((clock (cfg-actor-clock advanced actor)))
      (enqueue-over-states advanced (get-cfg-states cfg)
                           payload (lambda (_) clock))))

define (broadcast cfg payload) :=
  (enqueue-over-states cfg (get-cfg-states cfg)
                       payload (lambda (_) clock-empty))
\end{lstlisting}

In \lstinline{broadcast-from}, the sender's clock is advanced via \lstinline{cfg-tick}, and the resulting timestamp is attached to messages enqueued for every actor via \lstinline{enqueue-over-states}.  
The \lstinline{broadcast} function uses a zero vector clock for situations where timestamps are not required.

\subsection{Initialisation and progress conditions}

\textbf{TO DO: explain this definition of progress.}  
\textbf{TO DO: and the formalisation???}

We must create a configuration with initial states and clocks for each actor.  
The helper \lstinline{init-state-map} builds a map from a list of actor identifiers to the payload’s initial state,  
while \lstinline{init-clock-map} constructs a registry of zero vector clocks.  

\AthenaCode{lst:backbone-init-maps}{Initialising state and clock maps}
\begin{lstlisting}[language=Athena,style=athena]
define (init-state-map ops ids) :=
  (Payload.init-map ops ids)

define (init-clock-map ids) :=
  match ids {
    nil => registry-empty
  | (id :: rest) =>
      (FMap.update (pair id clock-empty) (init-clock-map rest))
  | _ => registry-empty
  }
\end{lstlisting}

\AthenaCode{lst:backbone-init-config}{Configuration initialisation and progress predicates}
\begin{lstlisting}[language=Athena,style=athena]
define (init ops ids) :=
  (Backbone.kfg (init-state-map ops ids) 
                (init-clock-map ids) 
                empty-mset)

define (reducible cfg) :=
  match cfg {
    (Backbone.kfg _ _ pending) =>
      (net-has-pending pending)
  }

define (deadlocked cfg) :=
  match cfg {
    (Backbone.kfg _ _ pending) =>
      (net-empty pending)
  }
\end{lstlisting}

The predicate \lstinline{reducible cfg} returns \lstinline{true} whenever there is at least one message awaiting delivery.  
Conversely, \lstinline{deadlocked cfg} holds when the message network is empty.

\section{Properties and verification}

Message passing correctness depends on a combination of vector‑clock ordering and semilattice semantics.  
\textbf{TO DO: fix the explanation — supervisor unclear.}

\section{Conclusion}

The \texttt{CRDTBackbone} module isolates the concerns of network replication from the algebraic behaviour of a CRDT.  
By combining vector‑clock metadata with a simple message scheduler, it implements the assumptions required by operation‑based CRDTs.  
When composed with a payload satisfying the semilattice axioms and monotonicity from Chapter~\ref{chap:payload},  
the backbone yields a replicated data type that converges regardless of network reorderings or duplicates.

\textbf{TO DO: always a use the word snippet when referencing code, I better replace that word with the ref to that snippet using the associated label.}
\endinput